\section{Introduction}

The McDonald-Kreitman (MK) test \citep{mcdonald1991}
is among the most widely used methods for detecting natural selection
at the molecular level.
By comparing the ratio of nonsynonymous to synonymous changes
within species (polymorphism) to that between species (divergence),
the MK test provides a direct test of the neutral theory prediction
that these ratios should be equal. Departures from this expectation,
often quantified by the neutrality index \citep[NI;][]{rand1996}
or the fraction of adaptive substitutions $\alpha$ \citep{smith2002},
have revealed pervasive positive selection across a wide range of taxa.
% Maybe one or two citations here? I think Matt's book has some examples
% we could pull from.

A key challenge for the standard MK test is that
slightly deleterious nonsynonymous mutations,
which contribute to polymorphism but rarely fix,
inflate the number of nonsynonymous polymorphism ($P_n$) relative to
nonsynonymous substitutions ($D_n$) and bias $\alpha$ downward
\citep{charlesworth1993, eyre-walker2006}. An extension of the
standard MK test, the asymptotic MK test addresses this limitation 
by estimating $\alpha$ as a function of derived allele frequency
and extrapolating to fixation,
effectively removing the contribution of weakly deleterious variants
\citep{messer2013, haller2017}.
Another extension, the polarized MK test, uses sequences from a second
outgroup species to assign fixed differences to specific lineages,
allowing lineage-specific inference of selection \citep{smith2002}.

Despite the centrality of MK-based methods to molecular evolution,
existing software implementations remain fragmented.
% I have been formatting the names of software using monospaced font.
% I like it aesthetically, but disregard as needed.
\texttt{DnaSP} \citep{librado2009} provides the standard test
through a Windows-only graphical interface.
The \texttt{iMKT} R package and web server \citep{murga-moreno2019}
implements several extensions
but requires pre-computed summary statistics
and is tightly coupled to specific databases.
The \texttt{asymptoticMK} web tool \citep{haller2017}
implements the asymptotic test
but does not accept raw sequence data.
% I added degenotate here since it is the original thing Scott and I
% wanted to test for the barnacles. I've also seen it used on a few papers,
% since it is part of snpArcher. When it works, it seems to work well, but
% it is dependent on having those outgroup variants which can get tricky on
% some cases. For example, for us the biggest setback was lifting the
% crenatus variants into the glandula genome to generate a single VCF with
% ingroup and outgroup sequences.
The \texttt{degenotate} software, part of the \texttt{snpArcher}
pipeline \cite{mirchandani2023}, can calculate both standard and 
imputed versions of the test; however it does so only using
genetic variants in VCF format and is thus dependent on
(and susceptible to) the liftover of outgroup sequences.
No existing tool provides a unified command-line interface
that takes aligned coding sequences as input
and supports the standard, polarized, and asymptotic MK tests
together with batch processing across thousands of genes
(\cref{tab:comparison}).

Here we present \mkado, a Python package that fills this gap.
\mkado takes codon-aligned FASTA files as input
and implements the standard, polarized, and asymptotic MK tests,
as well as related statistics,
with a consistent command-line interface.
It supports parallel batch processing of genome-scale datasets,
automatic multiple testing correction,
and built-in visualization
including volcano plots and asymptotic $\alpha$ curves.
\mkado is installable via \texttt{pip}
and is designed to integrate naturally
with modern Python-based bioinformatics workflows.

\begin{table}[!t]
  \centering
  \caption{Feature comparison of MK test software. Checkmarks indicate
    built-in support; dashes indicate the feature is absent.}
  \label{tab:comparison}
  \footnotesize
  \begin{tabular}{@{}lcccc@{}}
    \toprule
    Feature & \texttt{DnaSP} & \texttt{iMKT} & \begin{tabular}[c]{@{}c@{}}asymptotic\\[-3pt]MK\end{tabular} & \mkado \\
    \midrule
    Sequence input        & \checkmark & web only   & ---          & \checkmark \\
    Standard MK           & \checkmark & \checkmark & ---          & \checkmark \\
    Polarized MK          & ---        & ---        & ---          & \checkmark \\
    Asymptotic MK         & ---        & \checkmark & \checkmark   & \checkmark \\
    $\alpha_\text{TG}$    & ---        & ---        & ---          & \checkmark \\
    Frequency filtering   & ---        & \checkmark & \checkmark   & \checkmark \\
    CLI                   & ---        & ---        & ---          & \checkmark \\
    Python API            & ---        & ---        & ---          & \checkmark \\
    Batch processing      & ---        & ---        & ---          & \checkmark \\
    Parallel execution    & ---        & ---        & ---          & \checkmark \\
    Multiple testing corr.& ---        & ---        & ---          & \checkmark \\
    Visualization         & limited    & \checkmark & \checkmark   & \checkmark \\
    Open source           & ---        & \checkmark & \checkmark   & \checkmark \\
    Cross-platform        & ---        & \checkmark & \checkmark   & \checkmark \\
    \bottomrule
  \end{tabular}
\end{table}
% If we decide to add degenotate, it can do both a standard and the imputed MK test
% (which is the one from iMKT). It can also do DoS according to their docs.
% https://github.com/harvardinformatics/degenotate
% It takes a VCF input, has a CLI, is open source, and cross platform. I can take
% a more detailed look if we decide to add it.